\originalchapter{简并微扰与能级分裂}
\originalsection{简并子空间中的对角化}
若$H_0$有多维本征子空间$\mathcal D$, 非简并公式会出现零分母. 在一阶方程上用完整$E_0$本征空间投影$P$投到$\mathcal D$, 左边$P(H_0-E_0)=0$, 因而
$PWP\ket{\chi}=E^{(1)}\ket\chi$. 零阶态不应事先任取, 而要选$W$在简并子空间内的本征向量. 首先对角化一个较小的厄米矩阵, 才能决定一阶分裂与正确起始态.

\begin{example}
$H_0=\operatorname{diag}(E_0,E_0,E_3)$, $E_3\ne E_0$. 微扰在前两维为 $\begin{pmatrix}a&b\\b^*&a\end{pmatrix}$, $b=|b|\e^{\ii\varphi}\ne0$. 求一阶能量和零阶态.
\end{example}
\begin{solution}
特征方程是$(a-w)^2-|b|^2=0$, 所以$w_\pm=a\pm|b|$. 对应归一向量可取 $(\e^{\ii\varphi}\ket1\pm\ket2)/\sqrt2$, 能量为$E_0+\lambda(a\pm|b|)+O(\lambda^2)$. 与第3态的耦合先在子空间外修正中出现, 由非零能隙$E_0-E_3$控制. 当$E_3$也接近$E_0$, 应把它一并纳入近简并空间.
\end{solution}

若$PWP$仍有简并, 一阶没有完全区分态. 可以进一步研究二阶有效算符$PWQ(E_0-QH_0Q)^{-1}QWP$, $Q=I-P$. 此式要求补空间上对应逆算符存在, 不可在仍有零能隙时直接取逆. 一阶未分裂不等于扰动完全没有作用.

\originalsection{自旋轨道耦合}
在固定$n,\ell$的氢样子空间, 考虑$W=\xi(r)\boldsymbol L\cdot\boldsymbol S$. 若径向矩阵元已知, 用总角动量$\boldsymbol J=\boldsymbol L+\boldsymbol S$对角化角向部分最自然. $L^2,S^2,J^2,J_z$相容, 而$L_z,S_z$通常分别不守恒.

利用 $2\boldsymbol L\cdot\boldsymbol S=J^2-L^2-S^2$, 自旋$1/2$时, $j=\ell+1/2$的角向本征值为$\hbar^2\ell/2$, $j=\ell-1/2$为$-\hbar^2(\ell+1)/2$. 后者只在$\ell\ge1$存在. 所以一阶能量移为对应角向值乘$\langle\xi(r)\rangle_{n\ell}$. 这清楚分开了径向积分与角动量代数; 尚未求径向积分时, 不能把系数写成一个虚构数值.

氢精细结构还包含相对论动能修正与Darwin项, 并非仅有自旋轨道耦合. 完整结果的推导需这些有效Hamilton算符及其适用阶数; 本册此处完成角向分裂方法, 完整精细结构计算留作来源中的8.06第2章对应专题, 覆盖表标为尚未完整重编. 这样不会把“一个分裂项已算”误记成“整套精细结构已推导”.

\originalsection{弱场Zeeman效应}
弱外磁场相对于精细结构能隙很小, 仍可用$\ket{n,\ell,j,m_j}$为零阶态. 残余简并中,$H_Z$与$L^2,J_z$对易, 所以不同$\ell,m_j$不混合; 相同$\ell,m_j$的其他$j$分量由非零精细结构能隙分开. 忽略核磁矩, 电子磁扰动为$H_Z=\mu_B B(L_z+2S_z)/\hbar$, $\mu_B=e\hbar/(2m_e)>0$, 磁场沿$z$轴. 正号来自电子磁矩为负且能量为$-\boldsymbol\mu\cdot\boldsymbol B$.

在固定$j>0$的单个不可约角动量空间, 向量算符满足$[J_i,V_j]=\ii\hbar\epsilon_{ijk}V_k$. 标准投影性质给$P_j\boldsymbol S P_j=c_j\boldsymbol J$, 其中$c_j$与$m_j$无关. 两侧点乘$\boldsymbol J$并利用固定$J^2$, 就由$\boldsymbol S\cdot\boldsymbol J$确定$c_j$. 因此向量算符$\boldsymbol S$的对角期望沿$\boldsymbol J$的方向, 比例由$\boldsymbol S\cdot\boldsymbol J$决定. 取$2\boldsymbol S\cdot\boldsymbol J=J^2+S^2-L^2$, 得
\[
\langle S_z\rangle=\hbar m_j
\frac{j(j+1)+s(s+1)-\ell(\ell+1)}{2j(j+1)}.
\]
其比例关系是固定$j$空间中的向量投影性质, 可由角动量升降矩阵元验证; 这里引用此明确性质, 不把它视为对任意向量场成立. 因而$\Delta E_Z=\mu_BBg_jm_j$, 其中
$g_j=1+[j(j+1)+s(s+1)-\ell(\ell+1)]/[2j(j+1)]$, $s=1/2$. 若磁场超过精细结构尺度, 固定$j$近似失效, 需换用更适合的基并重新对角化. 使用公式必须先比较能隙.
